\section{Conclusiones}
El desarrollo de este caso de estudio para TECH TRANSITION permitió transformar un conjunto de requerimientos planos en un modelo de datos estructurado y relacional. Mediante la normalización, se separaron las entidades fundamentales (Clientes, Sucursales y Artículos) de los movimientos transaccionales (Cabecera y Detalle de Factura). La implementación en Oracle Database garantizó la integridad referencial a través de restricciones de llaves primarias (PK) y foráneas (FK), asegurando que no existan facturas huérfanas ni referencias a artículos inexistentes. Esta arquitectura robusta no solo optimiza el almacenamiento al eliminar redundancias, sino que establece las bases necesarias para futuras implementaciones de inteligencia de negocios o estructuras analíticas tipo Data Warehouse dentro del departamento de facturación.